% Mit twoside als weitere Option werden die Seitenränder angepasst, sodass beidseitig gedruckt werden kann
\documentclass[pdftex,12pt,a4paper]{report}

% kann optional mit "german" als Sprache gewählt werden
\usepackage[]{dbstmpl}



\global\arbeit{Bachelor Thesis}
%\global\arbeit{Masterarbeit}
%\global\arbeit{Seminararbeit}




% Hier die eigenen Daten eintragen
\global\titel{TITEL DER ARBEIT KANN AUCH EINE ZWEITE ZEILE HABEN}
\global\bearbeiter{BEARBEITERNAME}
\global\betreuer{BETREUERNAME}
\global\aufgabensteller{AUFGABENSTELLERNAME}
\global\abgabetermin{20.10.2017}
\global\ort{ORT}
\global\fach{Computer Science}

\begin{document}

% Deckblatt
\deckblatt

% Erklaerung fuer das Pruefungsamt
% Bei Seminararbeiten nicht noetig
\erklaerung

% Zusammenfassung
\begin{abstract}
Dieses Dokument dient als Muster für die Ausarbeitung einer \the\arbeit\
am Lehrstuhl für Datenbanksysteme am Institut für
Informatik der LMU München. Der Abstract sollte nicht mehr als 300 Wörter enthalten.
\end{abstract}

% Inhaltsverzeichnis (abhaengig von der Laenge des Dokuments)
\tableofcontents

% Hier beginnt der Hauptteil

\input{Introduction}









%% Ende Hauptteil

% Abbildungsverzeichnis (kann auch nach dem Inhaltsverzeichnis kommen)
%\listoffigures

% Tabellenverzeichnis (kann auch nach dem Inhaltsverzeichnis kommen)
%\listoftables

% Literaturverzeichnis
\bibliographystyle{plain} 
\bibliography{dbstmpl}         % verwendet dbstmpl.bib

\end{document}
